% Options for packages loaded elsewhere
\PassOptionsToPackage{unicode}{hyperref}
\PassOptionsToPackage{hyphens}{url}
\documentclass[
]{article}
\usepackage{ma2003b-notes}
\hypersetup{
  pdftitle={Multivariate Analysis - Lecture Notes},
  pdfauthor={MA2003B},
  hidelinks,
  pdfcreator={LaTeX via pandoc}}

\title{Multivariate Analysis - Lecture Notes}
\author{MA2003B}
\date{}

\begin{document}
\maketitle

\textbf{Multivariate Analysis Foundations}

Theoretical Principles, Probability Models, and Diagnostics

MA2003B --- Application of Multivariate Methods in Data Science

Tecnológico de Monterrey

\begin{center}\rule{0.5\linewidth}{0.5pt}\end{center}

\textbf{ }

\section{Course and Chapter Overview}

\textbf{ }

\subsection{Course Context}

Multivariate statistical methods extend classical single-variable
inference to systems where \(p \geq 2\) interconnected random variables
are observed simultaneously. In \textbf{MA2003B}, this module provides
the fundamental matrix algebra, probability theory, and diagnostic
toolset required for all downstream dimensionality reduction (PCA,
Factor Analysis) and supervised classification (Discriminant Analysis)
techniques.

\textbf{ }

\subsection{Learning Objectives}

By the end of this chapter, students will be able to:

\begin{itemize}
\item
  Formulate random vectors, joint probability density functions, and
  marginal/conditional expectations.
\item
  Compute and interpret sample mean vectors (\(\overline{\mathbf{x}}\)),
  sample covariance matrices (\(\mathbf{S}\)), and sample correlation
  matrices (\(\mathbf{R}\)).
\item
  Analyze the mathematical properties and geometry of the Multivariate
  Normal (MVN) distribution.
\item
  Construct constant-density contour ellipsoids and compute Mahalanobis
  generalized distances.
\item
  Identify missing data mechanisms (MCAR, MAR, MNAR) and apply modern
  imputation algorithms (KNN, MICE).
\item
  Detect multivariate outliers using Mahalanobis distances and Robust
  Minimum Covariance Determinant (MCD) estimators.
\item
  Derive Fisher \(z\)-transformations and construct confidence intervals
  for correlation parameters (\(\rho\)).
\item
  Apply multidimensional visualization techniques including scatter
  matrices, correlation heatmaps, and Andrews curves.
\end{itemize}

---

\textbf{ }

\section{2.1 Multivariate Distributions \& Random Vectors}

\textbf{ }

\subsection{Random Vectors}

Let
\(\mathbf{X} = \left\lbrack X_{1},X_{2},\ldots,X_{p} \right\rbrack^{T}\)
be a \(p\)-dimensional column random vector where each component
\(X_{j}\) is a real-valued random variable.

\textbf{ }

\subsection{Joint Cumulative Distribution Function (CDF)}

\[F\left( \mathbf{x} \right) = P\left( X_{1} \leq x_{1},X_{2} \leq x_{2},\ldots,X_{p} \leq x_{p} \right)\]

\textbf{ }

\subsection{Joint Probability Density Function (PDF)}

For continuous random vectors with density
\(f\left( \mathbf{x} \right) = f\left( x_{1},\ldots,x_{p} \right)\):
\[P\left( \mathbf{X} \in A \right) = \int\ldots\int_{A}f\left( x_{1},\ldots,x_{p} \right),dx_{1}\ldots dx_{p}\]

Marginal density for subset \(\mathbf{X}_{1}\):
\[f_{1}\left( \mathbf{x}_{1} \right) = \int_{\mathbb{R}^{p - k}}f\left( \mathbf{x}_{1},\mathbf{x}_{2} \right),d\mathbf{x}_{2}\]

Conditional density of \(\mathbf{X}_{1}\) given
\(\mathbf{X}_{2} = \mathbf{x}_{2}\):
\[f\left( \mathbf{x}_{1}~|~\mathbf{x}_{2} \right) = \frac{f\left( \mathbf{x}_{1},\mathbf{x}_{2} \right)}{f_{2}\left( \mathbf{x}_{2} \right)},\quad\text{ provided  }f_{2}\left( \mathbf{x}_{2} \right) > 0\]

---

\textbf{ }

\section{2.2 Mean Vectors and Covariance Matrices}

\textbf{ }

\subsection{Population Mean Vector}

\[\mathbf{\mu} = \mathbb{E}\left\lbrack \mathbf{X} \right\rbrack = \begin{bmatrix}
\mathbb{E}\left\lbrack X_{1} \right\rbrack \\
\mathbb{E}\left\lbrack X_{2} \right\rbrack \\
 \vdots \\
\mathbb{E}\left\lbrack X_{p} \right\rbrack
\end{bmatrix} = \begin{bmatrix}
\mu_{1} \\
\mu_{2} \\
 \vdots \\
\mu_{p}
\end{bmatrix}\]

\textbf{ }

\subsection{Population Covariance Matrix}

\[\mathbf{\Sigma} = \operatorname{Cov}(\mathbf{X}) = \mathbb{E}\left\lbrack \left( \mathbf{X} - \mathbf{\mu} \right)\left( \mathbf{X} - \mathbf{\mu} \right)^{T} \right\rbrack = \begin{bmatrix}
\sigma_{11} & \sigma_{12} & \ldots & \sigma_{1p} \\
\sigma_{21} & \sigma_{22} & \ldots & \sigma_{2p} \\
 \vdots & \vdots & \ddots & \vdots \\
\sigma_{p1} & \sigma_{p2} & \ldots & \sigma_{pp}
\end{bmatrix}\]

where \(\sigma_{jj} = \operatorname{Var}(X_{j}) = \sigma_{j}^{2}\) and
\(\sigma_{jk} = \operatorname{Cov}(X_{j},X_{k}) = \sigma_{kj}\).

\textbf{Info: } \textbf{Positive Semi-Definiteness:} For any non-zero
constant vector \(\mathbf{a} \in \mathbb{R}^{p}\),
\(\operatorname{Var}(\mathbf{a}^{T}\mathbf{X}) = \mathbf{a}^{T}\mathbf{\Sigma}\mathbf{a} \geq 0\).
Hence, all valid covariance matrices \(\mathbf{\Sigma}\) are symmetric
positive semi-definite (\(\mathbf{\Sigma} \succeq 0\)).

\textbf{ }

\subsection{Sample Statistics Matrix Operations}

Given a data matrix \(\mathbf{X} \in \mathbb{R}^{n \times p}\) with
\(n\) observations:

\begin{itemize}
\item
  \textbf{Sample Mean Vector:}
  \[\overline{\mathbf{x}} = \frac{1}{n}\mathbf{X}^{T}\mathbf{1}_{n} = \begin{bmatrix}
  {\overline{x}}_{1} & {\overline{x}}_{2} & \ldots & {\overline{x}}_{p}
  \end{bmatrix}^{T}\]
\item
  \textbf{Sample Covariance Matrix:}
  \[\mathbf{S} = \frac{1}{n - 1}\sum_{i = 1}^{n}\left( \mathbf{x}_{i} - \overline{\mathbf{x}} \right)\left( \mathbf{x}_{i} - \overline{\mathbf{x}} \right)^{T} = \frac{1}{n - 1}\mathbf{X}^{T}\left( \mathbf{I}_{n} - \frac{1}{n}\mathbf{1}_{n}\mathbf{1}_{n}^{T} \right)\mathbf{X}\]
\end{itemize}

---

\textbf{ }

\section{2.3 Correlations and Correlation Matrices}

\textbf{ }

\subsection{Population and Sample Correlation}

Let
\(\mathbf{D} = \operatorname{diag}(\sigma_{11},\sigma_{22},\ldots,\sigma_{pp})\)
be the diagonal \textbf{variance} matrix
(\(\sigma_{jj} = \operatorname{Var}(X_{j})\), not a standard deviation).
Its inverse square root
\(\mathbf{D}^{- \frac{1}{2}} = \operatorname{diag}(\frac{1}{\sigma_{1}},\ldots,\frac{1}{\sigma_{p}})\)
contains reciprocal standard deviations, which is what rescales
\(\mathbf{\Sigma}\) to a correlation matrix below.

\begin{itemize}
\item
  \textbf{Population Correlation Matrix:}
  \[\mathbf{P} = \mathbf{D}^{- \frac{1}{2}}\mathbf{\Sigma}\mathbf{D}^{- \frac{1}{2}} = \begin{bmatrix}
  1 & \rho_{12} & \ldots & \rho_{1p} \\
  \rho_{21} & 1 & \ldots & \rho_{2p} \\
   \vdots & \vdots & \ddots & \vdots \\
  \rho_{p1} & \rho_{p2} & \ldots & 1
  \end{bmatrix},\quad\rho_{jk} = \frac{\sigma_{jk}}{\sigma_{j}\sigma_{k}} \in \lbrack - 1,1\rbrack\]
\item
  \textbf{Sample Correlation Matrix:}
  \[\mathbf{R} = \mathbf{D}_{\mathbf{S}}^{- \frac{1}{2}}\mathbf{S}\mathbf{D}_{\mathbf{S}}^{- \frac{1}{2}},\quad r_{jk} = \frac{s_{jk}}{s_{j}s_{k}}\]
\end{itemize}

\textbf{ }

\subsection{Geometric Interpretation}

In sample space \(\mathbb{R}^{n}\), the sample correlation coefficient
\(r_{jk}\) equals the \textbf{cosine of the angle} \(\theta\) between
the mean-centered observation vectors
\(\mathbf{d}_{j} = \mathbf{x}_{j} - {\overline{x}}_{j}\mathbf{1}\) and
\(\mathbf{d}_{k} = \mathbf{x}_{k} - {\overline{x}}_{k}\mathbf{1}\):
\[r_{jk} = \cos(\theta) = \frac{\mathbf{d}_{j}^{T}\mathbf{d}_{k}}{|\mathbf{d}_{j}||\mathbf{d}_{k}|}\]

---

\textbf{ }

\section{2.4 The Multivariate Normal (MVN) Distribution}

\textbf{ }

\subsection{Probability Density Function}

A \(p\)-dimensional random vector \(\mathbf{X}\) follows a Multivariate
Normal distribution
\(\mathbf{X} \sim \mathcal{N}_{p\left( \mathbf{\mu},\mathbf{\Sigma} \right)}\)
(with positive definite \(\mathbf{\Sigma} \succ 0\)) if its joint PDF
is:

\[f\left( \mathbf{x} \right) = \frac{1}{(2\pi)^{\frac{p}{2}}|\mathbf{\Sigma}|^{\frac{1}{2}}}\exp( - \frac{1}{2}\left( \mathbf{x} - \mathbf{\mu} \right)^{T}\mathbf{\Sigma}^{- 1}\left( \mathbf{x} - \mathbf{\mu} \right))\]

\textbf{ }

\subsection{Key Properties of MVN}

\begin{enumerate}
\item
  \textbf{Linear Combinations:} If
  \(\mathbf{X} \sim \mathcal{N}_{p\left( \mathbf{\mu},\mathbf{\Sigma} \right)}\)
  and \(\mathbf{A} \in \mathbb{R}^{q \times p}\), then
  \(\mathbf{A}\mathbf{X} + \mathbf{b} \sim \mathcal{N}_{q\left( \mathbf{A}\mathbf{\mu} + \mathbf{b},\mathbf{A}\mathbf{\Sigma}\mathbf{A}^{T} \right)}\).
\item
  \textbf{Marginals and Conditionals:} All marginal and conditional
  distributions of an MVN vector are themselves Multivariate Normal.
\item
  \textbf{Zero Covariance Implies Independence:} For MVN distributions,
  \(\operatorname{Cov}(X_{j},X_{k}) = 0 \Longleftrightarrow X_{j} \perp X_{k}\).
\end{enumerate}

\textbf{ }

\subsection{Geometry of Contour Ellipsoids \& Mahalanobis Distance}

The exponent in the MVN density defines the \textbf{squared Mahalanobis
distance}:
\[D^{2}\left( \mathbf{x},\mathbf{\mu} \right) = \left( \mathbf{x} - \mathbf{\mu} \right)^{T}\mathbf{\Sigma}^{- 1}\left( \mathbf{x} - \mathbf{\mu} \right)\]

\begin{itemize}
\item
  The surfaces of constant probability density are hyper-ellipsoids
  centered at \(\mathbf{\mu}\):
  \[\{\mathbf{x} \in \mathbb{R}^{p}:\left( \mathbf{x} - \mathbf{\mu} \right)^{T}\mathbf{\Sigma}^{- 1}\left( \mathbf{x} - \mathbf{\mu} \right) = c^{2}\}\]
\item
  The principal axes of the ellipsoid align with the eigenvectors
  \(\mathbf{v}_{1},\ldots,\mathbf{v}_{p}\) of \(\mathbf{\Sigma}\), with
  half-lengths proportional to \(\pm c\sqrt{\lambda_{j}}\).
\item
  Under multivariate normality:
  \[D^{2}\left( \mathbf{X},\mathbf{\mu} \right) \sim \chi_{p}^{2}\]
\end{itemize}

---

\textbf{ }

\section{2.5 Lost, Null, and Incorrect Values (Data Quality \&
Imputation)}

\textbf{ }

\subsection{Missing Data Mechanisms}

\begin{itemize}
\item
  \textbf{Missing Completely at Random (MCAR):} The probability of
  missingness is independent of both observed and unobserved data:
  \(P\left( M~|~Y_{\text{obs}},Y_{\text{mis}} \right) = P(M)\).
\item
  \textbf{Missing at Random (MAR):} Missingness depends on observed
  variables but not unobserved values:
  \(P\left( M~|~Y_{\text{obs}},Y_{\text{mis}} \right) = P\left( M~|~Y_{\text{obs}} \right)\).
\item
  \textbf{Missing Not at Random (MNAR):} Missingness depends directly on
  the unobserved value itself.
\end{itemize}

\textbf{ }

\subsection{Imputation Strategies}

\begin{enumerate}
\item
  \textbf{Mean/Median Imputation:} Simple but distorts variance and
  attenuates covariance structures.
\item
  \textbf{K-Nearest Neighbors (KNN) Imputation:} Replaces missing values
  with the weighted average of the \(K\) closest complete records using
  Euclidean/Gower distance.
\item
  \textbf{Iterative Imputer (MICE - Multivariate Imputation by Chained
  Equations):} Models each feature with missing values as a function of
  all other features in a round-robin Gibbs sampling scheme.
\end{enumerate}

---

\textbf{ }

\section{2.6 Multivariate Aberrant Data (Outlier Detection)}

\textbf{ }

\subsection{Univariate vs. Multivariate Outliers}

An observation may have perfectly normal marginal values on each
individual coordinate \(X_{j}\) while being a severe multivariate
outlier because it violates the joint correlation structure
\(\mathbf{\Sigma}\).

\textbf{ }

\subsection{Diagnostic Metrics}

\begin{itemize}
\item
  \textbf{Sample Mahalanobis Squared Distance:}
  \[D_{i}^{2} = \left( \mathbf{x}_{i} - \overline{\mathbf{x}} \right)^{T}\mathbf{S}^{- 1}\left( \mathbf{x}_{i} - \overline{\mathbf{x}} \right)\]
\item
  \textbf{Chi-Square Outlier Rule:} Flag observation \(i\) as an outlier
  if \(D_{i}^{2} > \chi_{p,1 - \alpha}^{2}\) (typically
  \(\alpha = 0.001\)).
\item
  \textbf{Robust Covariance Estimation (MCD):} Classical
  \(\overline{\mathbf{x}}\) and \(\mathbf{S}\) are sensitive to outliers
  (breakdown point \(= \frac{1}{n}\)). The \textbf{Minimum Covariance
  Determinant (MCD)} finds the subset of \(h \leq n\) observations whose
  sample covariance matrix has the smallest determinant, providing
  robust location \({\hat{\mathbf{\mu}}}_{\text{MCD}}\) and scatter
  \({\hat{\mathbf{\Sigma}}}_{\text{MCD}}\).
\end{itemize}

---

\textbf{ }

\section{2.7 Sample Correlations: Fisher and Ruben Confidence Intervals}

\textbf{ }

\subsection{\texorpdfstring{Fisher's
\(z\)-Transformation}{Fisher's z-Transformation}}

The sampling distribution of Pearson's \(r\) is highly skewed when
\(|\rho| > 0\). Fisher's transformation normalizes the distribution:

\[z = \frac{1}{2}\ln(\frac{1 + r}{1 - r}) = \operatorname{arctanh}(r)\]

Asymptotically:
\[z \sim^{\text{approx}}\mathcal{N}(\frac{1}{2}\ln(\frac{1 + \rho}{1 - \rho}),\frac{1}{n - 3})\]

\(100(1 - \alpha)\%\) Confidence Interval for
\(\zeta = \operatorname{arctanh}(\rho)\):
\[\left\lbrack z_{L},z_{U} \right\rbrack = z \pm z_{1 - \frac{\alpha}{2}}\frac{1}{\sqrt{n - 3}}\]

Back-transforming to correlation scale:
\[\rho_{L} = \tanh(z_{L}) = \frac{\exp(2z_{L}) - 1}{\exp(2z_{L}) + 1},\quad\quad\rho_{U} = \tanh(z_{U}) = \frac{\exp(2z_{U}) - 1}{\exp(2z_{U}) + 1}\]

---

\textbf{ }

\section{2.8 Multivariate Descriptive Analytics \& Visualization}

\textbf{ }

\subsection{Multidimensional Plotting Techniques}

\begin{itemize}
\item
  \textbf{Scatter Plot Matrix (SPLOM):} Visualizes all
  \(\frac{p(p - 1)}{2}\) pairwise bivariate relationships simultaneously
  with diagonal univariate kernel density estimates.
\item
  \textbf{Correlation Heatmaps:} Displays \(\mathbf{R}\) with color
  intensity gradients and optional hierarchical clustering dendrograms.
\item
  \textbf{Andrews Curves:} Maps each \(p\)-dimensional observation
  \(\mathbf{x}_{i} = \left\lbrack x_{i1},\ldots,x_{ip} \right\rbrack^{T}\)
  to a continuous trigonometric function over
  \(t \in \lbrack - \pi,\pi\rbrack\):
  \[f_{\mathbf{x}_{i}}(t) = \frac{x_{i1}}{\sqrt{2}} + x_{i2}\sin(t) + x_{i3}\cos(t) + x_{i4}\sin(2t) + x_{i5}\cos(2t) + \ldots\]
  Similar multivariate observations appear as closely clustered curves
  in function space.
\end{itemize}

---

\textbf{ }

\section{Practical Python Implementation}

\begin{Shaded}
\begin{Highlighting}[]
\ImportTok{import}\NormalTok{ numpy }\ImportTok{as}\NormalTok{ np}
\ImportTok{import}\NormalTok{ pandas }\ImportTok{as}\NormalTok{ pd}
\ImportTok{from}\NormalTok{ scipy }\ImportTok{import}\NormalTok{ stats}
\ImportTok{from}\NormalTok{ sklearn.impute }\ImportTok{import}\NormalTok{ KNNImputer}
\ImportTok{from}\NormalTok{ sklearn.covariance }\ImportTok{import}\NormalTok{ MinCovDet}

\CommentTok{\# 1. Imputation of Missing Values (KNN)}
\NormalTok{imputer }\OperatorTok{=}\NormalTok{ KNNImputer(n\_neighbors}\OperatorTok{=}\DecValTok{5}\NormalTok{)}
\NormalTok{X\_imputed }\OperatorTok{=}\NormalTok{ pd.DataFrame(imputer.fit\_transform(df.iloc[:, }\DecValTok{1}\NormalTok{:]), columns}\OperatorTok{=}\NormalTok{df.columns[}\DecValTok{1}\NormalTok{:])}

\CommentTok{\# 2. Sample Mean Vector \& Covariance Matrix}
\NormalTok{mean\_vec }\OperatorTok{=}\NormalTok{ X\_imputed.mean()}
\NormalTok{cov\_mat }\OperatorTok{=}\NormalTok{ X\_imputed.cov()}
\NormalTok{corr\_mat }\OperatorTok{=}\NormalTok{ X\_imputed.corr()}

\CommentTok{\# 3. Mahalanobis Distance Outlier Detection}
\NormalTok{inv\_cov }\OperatorTok{=}\NormalTok{ np.linalg.inv(cov\_mat)}
\NormalTok{diff }\OperatorTok{=}\NormalTok{ X\_imputed }\OperatorTok{{-}}\NormalTok{ mean\_vec}
\NormalTok{mahal\_d2 }\OperatorTok{=}\NormalTok{ np.}\BuiltInTok{sum}\NormalTok{(diff.dot(inv\_cov) }\OperatorTok{*}\NormalTok{ diff, axis}\OperatorTok{=}\DecValTok{1}\NormalTok{)}
\NormalTok{cutoff }\OperatorTok{=}\NormalTok{ stats.chi2.ppf(}\FloatTok{0.999}\NormalTok{, df}\OperatorTok{=}\NormalTok{X\_imputed.shape[}\DecValTok{1}\NormalTok{])}
\NormalTok{outliers }\OperatorTok{=}\NormalTok{ X\_imputed[mahal\_d2 }\OperatorTok{\textgreater{}}\NormalTok{ cutoff]}
\BuiltInTok{print}\NormalTok{(}\SpecialStringTok{f"Detected }\SpecialCharTok{\{}\BuiltInTok{len}\NormalTok{(outliers)}\SpecialCharTok{\}}\SpecialStringTok{ multivariate outliers (cutoff = }\SpecialCharTok{\{}\NormalTok{cutoff}\SpecialCharTok{:.2f\}}\SpecialStringTok{)"}\NormalTok{)}

\CommentTok{\# 4. Fisher z Confidence Interval for Correlation}
\NormalTok{r }\OperatorTok{=}\NormalTok{ corr\_mat.loc[}\StringTok{"pm25"}\NormalTok{, }\StringTok{"pm10"}\NormalTok{]}
\NormalTok{n }\OperatorTok{=} \BuiltInTok{len}\NormalTok{(X\_imputed)}
\NormalTok{z }\OperatorTok{=}\NormalTok{ np.arctanh(r)}
\NormalTok{se }\OperatorTok{=} \DecValTok{1} \OperatorTok{/}\NormalTok{ np.sqrt(n }\OperatorTok{{-}} \DecValTok{3}\NormalTok{)}
\NormalTok{ci\_z }\OperatorTok{=}\NormalTok{ (z }\OperatorTok{{-}} \FloatTok{1.96} \OperatorTok{*}\NormalTok{ se, z }\OperatorTok{+} \FloatTok{1.96} \OperatorTok{*}\NormalTok{ se)}
\NormalTok{ci\_r }\OperatorTok{=}\NormalTok{ (np.tanh(ci\_z[}\DecValTok{0}\NormalTok{]), np.tanh(ci\_z[}\DecValTok{1}\NormalTok{]))}
\BuiltInTok{print}\NormalTok{(}\SpecialStringTok{f"PM2.5 vs PM10: r = }\SpecialCharTok{\{}\NormalTok{r}\SpecialCharTok{:.4f\}}\SpecialStringTok{, 95\% CI = [}\SpecialCharTok{\{}\NormalTok{ci\_r[}\DecValTok{0}\NormalTok{]}\SpecialCharTok{:.4f\}}\SpecialStringTok{, }\SpecialCharTok{\{}\NormalTok{ci\_r[}\DecValTok{1}\NormalTok{]}\SpecialCharTok{:.4f\}}\SpecialStringTok{]"}\NormalTok{)}
\end{Highlighting}
\end{Shaded}


\end{document}

